\chapter{Prenatal Cell-Free DNA (cfDNA)}

\section*{What Is This Test?}
Prenatal cell-free DNA (cfDNA) screening, also known as noninvasive prenatal testing (NIPT) or noninvasive prenatal screening (NIPS), is a method of screening for fetal chromosomal aneuploidies by analyzing cell-free fetal DNA (cffDNA) circulating in maternal blood. During pregnancy, approximately 10--20\% of cell-free DNA in maternal plasma originates from the placenta (specifically from apoptosis of cytotrophoblast cells) and reflects the fetal chromosomal complement. cfDNA screening analyzes this mixture of maternal and placental cell-free DNA to detect trisomy 21 (Down syndrome, sensitivity >99\%, false-positive rate <0.1\%), trisomy 18 (Edwards syndrome, sensitivity 97\%, false-positive rate <0.1\%), trisomy 13 (Patau syndrome, sensitivity 90--95\%), and sex chromosome aneuploidies (monosomy X [Turner], 47,XXX, 47,XXY [Klinefelter], 47,XYY) \cite{gil2017analysis,norton2015cellfree}. Some expanded cfDNA panels also screen for microdeletion syndromes (22q11.2 deletion/DiGeorge syndrome, 5p deletion/Cri-du-chat, 15q11.2/Angelman/Prader-Willi) and rare autosomal trisomies. cfDNA screening is performed as early as 9--10 weeks of gestation (when the fetal fraction of placental DNA is $\geq$4\%) and is currently recommended by the American College of Obstetricians and Gynecologists (ACOG) as the most sensitive and specific screening test for fetal trisomies, offered to all pregnant women regardless of maternal age or baseline risk \cite{acog2020screening}.

\section*{Alternative Names}
NIPT; NIPS; Noninvasive Prenatal Testing/Screening; Cell-Free Fetal DNA; cffDNA; cfDNA Screening; Maternal Plasma DNA Testing; Prenatal Aneuploidy Screening

\section*{Principle}
Maternal blood is collected into a specialized cfDNA collection tube (Streck or Roche cell-free DNA tube) that contains a preservative to stabilize nucleated blood cells and prevent maternal DNA release (and thus dilution of fetal cfDNA). Plasma is separated within 6--7 days. Cell-free DNA is extracted and analyzed by massively parallel shotgun sequencing (MPSS), which reads and counts millions of short DNA fragments from the entire genome. Each chromosome's representation is quantified. In a euploid pregnancy, the ratio of chromosome-specific sequence reads is proportional to the chromosome copy number. An overrepresentation of sequences from chromosome 21 (Z-score >3) indicates trisomy 21. An alternative method, targeted SNP-based sequencing (Panorama), uses multiplex PCR amplification of specific SNPs, followed by NGS and bioinformatics analysis using the Next-Generation Aneuploidy Test Using SNPs (NATUS) algorithm. A third method, microarray-based comparative genomic hybridization (CGH-array), is used in some platforms (Verifi, Illumina). The fetal fraction (FF) --- the percentage of cell-free DNA in maternal plasma that is of placental/fetal origin --- is calculated; an FF <4\% (insufficient fetal DNA) is reported as ``no call'' or ``test failure'' and requires a repeat draw.

\section*{History}
The discovery of cell-free fetal DNA in maternal plasma by Dennis Lo and colleagues at Oxford University in 1997 revolutionized prenatal screening \cite{lo1997presence}. Initially, cffDNA was used for fetal sex determination and fetal RhD genotyping. The application of massively parallel sequencing for aneuploidy detection was pioneered by Fan, Chiu, and Lo in 2008, followed by multiple validation studies (Chiu et al., 2008; Palomaki et al., 2011; Bianchi et al., 2012) \cite{palomaki2011dnasequencing,bianchi2018sequencing}. NIPT became clinically available in the United States in late 2011 (MaterniT21 by Sequenom, Verifi by Illumina, Panorama by Natera, Harmony by Ariosa). ACOG and the Society for Maternal-Fetal Medicine (SMFM) initially recommended NIPT for high-risk pregnancies (maternal age $\geq$35, abnormal ultrasound, abnormal serum screening) in 2012, but by 2020, NIPT is recommended for all pregnant women, regardless of risk. The FDA has not directly approved NIPT but classifies it as a laboratory-developed test (LDT). The 2019--2021 ACOG Practice Bulletins established NIPT as a first-tier screening test.

\section*{Why This Test Is Ordered}
Screening for fetal aneuploidy (trisomy 21, 18, 13) and sex chromosome aneuploidy in all pregnant women, per ACOG 2020 recommendation. NIPT is the preferred first-trimester screening test, particularly for women of advanced maternal age. Screening for fetal sex and sex chromosome aneuploidy (optional). Screening for microdeletion syndromes (22q11.2, 5p, 15q) when expanded NIPT is requested, though positive predictive value (PPV) is significantly lower for microdeletions. Detection of fetal sex for X-linked disorders (optional). In RhD-negative women, cfDNA can detect the fetal RHD gene, guiding Rh immunoglobulin (RhoGAM) administration.

\section*{Primary vs Secondary Test}
This is a primary screening test (NOT a diagnostic test). A positive cfDNA result MUST be confirmed by diagnostic testing (amniocentesis or CVS) before any irreversible clinical action is taken.

\section*{How This Test Is Performed}
Maternal blood sample (10--20 mL) is collected in a specialized cfDNA collection tube between 10 weeks and term. The specimen must be maintained at room temperature and processed within 6--7 days.

\section*{What Preparation Is Needed}
No special preparation is required, and no fasting is necessary. The blood sample may be drawn at any time of day. Because the test relies on a sufficient proportion of circulating placental (fetal) DNA, it should not be performed before 10 weeks of gestation, when the fetal fraction may be too low for reliable analysis. In RhD-negative women undergoing fetal \textit{RHD} genotyping, no preparation is needed, and prior administration of Rh immune globulin (RhoGAM) does not interfere with the analysis.

\section*{Contraindications and Precautions}
There are no absolute contraindications to the blood draw. The test is not validated for all multiple gestations: most platforms can assess twin pregnancies, but higher-order multiples and some twin situations (e.g., egg donation, surrogacy) should be discussed with the laboratory. Precautions: cfDNA screening is a screening test, not a diagnostic test, and a positive result MUST be confirmed by chorionic villus sampling (CVS) or amniocentesis before any irreversible action is taken. Factors that lower the fetal fraction (maternal obesity with BMI >30--35, early gestational age, low-molecular-weight heparin therapy, maternal malignancy, autoimmune disease) increase the chance of a ``no call'' result or a false-negative result. Confined placental mosaicism and maternal copy-number variation or maternal sex chromosome mosaicism can produce discordant (false-positive or false-negative) results.

\section*{Risks}
Minimal risk. The test requires only a standard maternal venipuncture and carries no risk of miscarriage or injury to the fetus, unlike invasive diagnostic procedures (chorionic villus sampling or amniocentesis). Slight pain, bruising, or, rarely, hematoma at the venipuncture site.

\section*{What Happens After the Test}
Results are typically available within 5--10 days after the laboratory receives the sample. A negative (low-risk) result does not eliminate the need for second-trimester maternal serum alpha-fetoprotein (MSAFP) screening and the 18--22 week anatomy ultrasound, which detect open neural tube defects and structural anomalies not assessed by cfDNA. A positive (high-risk) result is followed by genetic counseling and confirmatory diagnostic testing (CVS or amniocentesis) before any irreversible decision is made.

\section*{Understanding Results}

\subsection*{Negative (Low Risk/Low Probability)}
The test result is ``low risk'' or ``negative'' for the chromosomes tested. The probability of trisomy 21, 18, and 13 is very low. However, a negative result does NOT guarantee a chromosomally normal fetus. False negatives occur in 0.1--0.5\% for trisomy 21. Standard prenatal care continues per ACOG guidelines, including second-trimester maternal serum alpha-fetoprotein (MSAFP) screening for neural tube defects (which are NOT detected by cfDNA) and anatomy ultrasound at 18--22 weeks.

\subsection*{Positive (High Risk/High Probability)}
The test result indicates a high probability of fetal chromosomal aneuploidy for a specific chromosome. Positive predictive value (PPV) varies: for trisomy 21, PPV is approximately 90--95\% in high-risk populations, 80--90\% in average-risk populations. For trisomy 18: PPV approximately 50--80\%. For trisomy 13: PPV approximately 30--60\%. For monosomy X (Turner): PPV approximately 25--50\% (lower due to confined placental mosaicism and maternal sex chromosome mosaicism/age-related X chromosome loss). Microdeletion (22q11.2): PPV 20--50\%. A positive result MUST be confirmed by diagnostic testing (chorionic villus sampling [CVS] at 10--13 weeks or amniocentesis at 15--20 weeks) with karyotype or chromosomal microarray analysis (CMA).

\subsection*{No Call/Test Failure}
The fetal fraction (FF) is <4\%, or the test result is indeterminate. Occurs in 1--5\% of samples. Causes: early gestational age (<10 weeks), maternal obesity (BMI >30--35 decreases fetal fraction), aneuploidy (trisomy 13, 18, and triploidy are associated with low FF), maternal autoimmune disease, medications (low-molecular-weight heparin), maternal malignancy (tumor DNA dilutes fetal fraction). Repeat draw recommended. If repeat fails, diagnostic testing (CVS or amniocentesis) should be offered because low fetal fraction is associated with an increased risk of fetal aneuploidy.

\section*{Normal Results}
Low risk/negative for trisomy 21, 18, and 13, and sex chromosome aneuploidy. Fetal fraction $\geq$4\% (sufficient for reliable analysis). cfDNA screening does NOT detect neural tube defects; MSAFP screening at 15--22 weeks and anatomy ultrasound at 18--22 weeks are still required.

\section*{Follow-up Tests}
Diagnostic testing (chorionic villus sampling [CVS] or amniocentesis) with karyotype or chromosomal microarray (CMA) to confirm a positive cfDNA result. Second-trimester MSAFP screening (at 15--22 weeks) for open neural tube defects (anencephaly, open spina bifida), which are NOT detected by cfDNA. Fetal anatomy ultrasound at 18--22 weeks. Genetic counseling to discuss results, implications, and follow-up options.

\section*{References}
\bibliographystyle{unsrtnat}
\bibliography{references}

\clearpage
